\color{unidarkgreen}{\subsection{Автоматическое ускорение при включении выключателя (АУ)}}
\color{black}

\begin{enumerate}
%%%%%%%%%%%%%%%%%%%%%%%%%%%%%%%%%%%%%%%%%%%%%%%%%%%%%%%%%%%%%%%%%%%%%%  ОБЩИЕ СВЕДЕНИЯ %%%%%%%%%%%%%%%%%%%%%%%%%%%%%%%%%%%%%%%%%%%%%%%%%%

\item
В устройстве реализована логика автоматического ускорения (АУ) ДЗ, ТНЗНП и АМТЗ при включении выключателя, которая вводится программным ключом \textbf{«АУ»}.

\item
Автоматическое ускорение вводится на время, задаваемое уставкой \textbf{«Тв АУ»}.

\item
Имеется возможность задать контроль отсутствия напряжения на линии при вводе АУ с помощью программного ключа \textbf{«Контр. отс. U на ЛЭП»}. При установке измерительного трансформатора напряжения на шинах в качестве источника информации о напряжении на линии может использоваться переменная Uлэп (способ её настройки приведен в пункте \ref{func:nkn}) или внешний сигнал от ИО минимального напряжения. Способ контроля отсутствия напряжения на линии выбирается с помощью программного переключателя \textbf{«Контр. U ЛЭП для АУ»}.
Если измерительный трансформатор напряжения установлен на линии (определяется положением переключателя \textbf{«Место ТН»}), то контролируется отсутствие симметричного напряжения. В этому случае для ввода АУ напряжение прямой последовательности на линии не должно превышать уставку \textbf{«U1 макс АУ»}, а напряжения обратной и нулевой последовательности должны быть больше уставок \textbf{«U2 макс АУ»} и \textbf{«3U0 макс АУ»} соответственно.

\item
При необходимости можно ввести контроль наличия симметричного напряжения на смежном элементе (шинах), который задается программным ключом \textbf{«Контр. нал. U на шинах»}. Наличие напряжения на энергообъекте 1 определяется по напряжению Uэ1 (настройка переменной Uэ1 также приведена в пункте \ref{func:nkn}), которое для ввода АУ должно превышать уставку \textbf{«Uнал шин Э1»}. В случае присоединения линии к системе через два выключателя наличие напряжения на энергообъекте Э2 определяется по уставке \textbf{«Uнал шин Э2»}.

\item
Условия по напряжению контролируются только в момент, предшествующий включению.

\item
Выдержка времени срабатывания ускоряемой ступени ДЗ регулируется уставкой \textbf{«Тср АУ ДЗ»}, ступени ТНЗНП -- \textbf{«Тср АУ ТНЗНП»}, вторых ступеней АМТЗ -- \textbf{«Тср АУ АМТЗ»}.

\item
АУ АМТЗ вводится программным ключом \textbf{«АУ АМТЗ»}.

\item
Функциональная схема логики АУ приведена на рисунке \ref{pict_au:AU}. 

\vspace{2mm}
\begin{figure}[H]
\centering
\includegraphics[width=0.9\textwidth,height=0.9\textheight,keepaspectratio]{\fbpath/03.01 ЛВ/561.1311 АУ/img/AU.pdf}
\captionof{figure}{Функциональная схема АУ}\label{pict_au:AU}
\end{figure}

\item
Уставки АУ представлены в таблице \ref{app_set:au}.

\end{enumerate}